\subsection{The splitting, Hopf and Riedler families}
\label{sec:quartic-family-completions}

We apply the common theorem to the splitting construction, the Hopf
pairing and Riedler's pullbacks. The splitting and Riedler families
admit complete variational classifications. The balanced Hopf section
also illustrates intrinsic comparison inside a Simons cone, and the
degree-nine Riedler potential supplies an additional Hardy improvement.

\subsubsection{The OT--FKM splitting criterion}

For the data \eqref{bh:quartic-data}--\eqref{bh:centered}, the exact
minimizing threshold is as follows.

\begin{theorem}\label{bh:main}
The hypersurface $C=\{F=0\}\subset\R^{2\ell}$ is minimal on its
regular part, and
\[
 \Sing C=(\R^\ell\times\{0\})\cup(\{0\}\times\R^\ell).
\]
The oriented boundary $\partial\llbracket\{F<0\}\rrbracket$ is
area minimizing if and only if $\ell\geq10$. In all these minimizing
cases it is strictly area minimizing in the sense of
\cref{def:four-variational-notions}. The polynomial potentials and
strict-stability bound of \cref{thm:quartic-table-existence} apply,
with degrees ten for $\ell=10$ and eight for $\ell\geq12$.
For the remaining
admissible dimensions $\ell=4,6,8$, the regular part of $C$ is unstable.
\end{theorem}

\paragraph{Closure and singularities.}

For $x\ne0$, the vectors \(x,A_1x,\ldots,A_{m-1}x\)
are orthogonal and have length $|x|$. Consequently $0\leq Q\leq T$.
Direct differentiation gives, in the variables $I=(u,v,Q)$,
\begin{equation}\label{bh:original-table}
 (\Gamma(I_i,I_j))=
 \begin{pmatrix}4u&0&4Q\\0&4v&4Q\\4Q&4Q&4RQ\end{pmatrix},
 \qquad
 \Delta I=(2\ell,2\ell,2mR).
\end{equation}
For example,
$\nabla_yQ=2\sum_{i=0}^{m-1}\ip{A_i x}{y}A_i x$, with $A_0=\Id$,
so $\norm{\nabla_yQ}^2=4uQ$; differentiating in $x$ gives $4vQ$.
By \cref{prop:MR-generator-test},
\begin{equation}\label{bh:algebra}
 \cA=\R[u,v,Q]=\R[R,u-v,F]
\end{equation}
is a graded Laplacian algebra, with generator degrees $(2,2,4)$.
The determinant of its gradient matrix is $64RQ(T-Q)$, which is
positive at points with $u,v>0$ and $0<Q<T$. These points exist by
choosing the components of $y$ both in and perpendicular to the
span above. Thus the three generators are algebraically independent;
the spherical quotient dimension is two. The separately harmonic
normal form of \cref{prop:two-block-normal} is
$P=Q-(m/\ell)uv$; the minimal polynomial $F=Q-cuv$ differs from
it by a radial term in the two blocks.

With
$p_0=u-v$ and $p_{i+1}=2\langle A_i x,y\rangle$ for
$A_0=\Id$, the matrices in \cref{rem:split-composed-foliation}
form a symmetric Clifford system, and
\[
 F=\frac14\left(\sum_{i=1}^{m}p_i^2
                    +cp_0^2-cR^2\right).
\]
Its blocks $E_0=p_0^2$, $E_1=4Q$, and $E_2=4(T-Q)$
have weights $(1,m,\ell-m)$. The substitutions in
\eqref{bh:centered} therefore give \eqref{bh:reduced-table} with
$h=\ell$ and $N=2\ell$.
The calculation after \eqref{bh:reduced-table} recovers
$\cA_R(F)=\R[R,T,F]$, of degrees $(2,4,4)$. Independence follows
since $(u,v,Q)\mapsto(u+v,uv,Q-cuv)$ has Jacobian $u-v$.
Both algebras have spherical quotient dimension two, but
\eqref{bh:algebra} also retains the sign of $u-v$.

The geometry follows from \cref{lem:quartic-table-geometry}. For
$\ell\ge10$, evenness of $\ell$ and
\cref{thm:quartic-table-existence} give the variational conclusions.
The remaining dimensions are excluded by the following local
model and angular calculation.

Fix $x_0\ne0$, and write
$W_{x_0}=\operatorname{span}\{x_0,A_1x_0,\ldots,A_{m-1}x_0\}$.
The first nonzero term of $F$ at $(x_0,0)$ is
\begin{equation}\label{bh:transverse}
 \norm{x_0}^2\big((1-c)\norm{y_\parallel}^2
                         -c\norm{y_\perp}^2\big),
 \quad y_\parallel\in W_{x_0},\quad y_\perp\in W_{x_0}^\perp.
\end{equation}
The tangent cone is the product of $\R^\ell$ with this nonlinear
Lawson cone. Its transverse link is
\[
 \Sph^{m-1}(\sqrt c)\times
 \Sph^{\ell-m-1}(\sqrt{1-c})\subset\Sph^{\ell-1}.
\]
Both multiplicities are positive by \eqref{bh:range}. Smooth local
orthonormal frames for $W_x$ and $W_x^\perp$ identify a neighborhood
of the stratum with a smoothly varying family of these quadratic cones.
The same description holds on the other coordinate subspace. 

\begin{remark}[Clifford coordinates of the OT--FKM splitting]
\label{rem:split-composed-foliation}
Put $A_0=\Id$. The symmetric Clifford system
\[
 P_0=\begin{pmatrix}\Id&0\\0&-\Id\end{pmatrix},\qquad
 P_{i+1}=\begin{pmatrix}0&A_i^T\\A_i&0\end{pmatrix}
 \quad(0\leq i\leq m-1)
\]
has coordinates $p_0=u-v$ and $p_{i+1}=2\langle A_i x,y\rangle$.
The common geometric interpretation is given in
\cref{sec:Clifford-composition-background}.
\end{remark}


\subsubsection{The unstable OT--FKM splitting dimensions}

For $\ell=4,6$, the transverse Lawson cone in \eqref{bh:transverse}
is unstable: its link has dimension $\ell-2$ and constant squared
curvature $\ell-2$, giving the Hardy margin
\begin{equation}\label{bh:Lawson-margin}
 \frac{(\ell-3)^2}{4}-(\ell-2)<0.
\end{equation}
\Cref{lem:tangent-instability-transfer} transfers a negative
variation from this transverse cone to $C_{\rm reg}$, using the
local model \eqref{bh:transverse}.

For $\ell=8$, the transverse Lawson cone can be area minimizing;
instability follows instead from the angular calculation below.

\begin{lemma}\label{bh:curvature}
On the regular part of the quartic cone,
\begin{equation}\label{bh:AC}
 \norm{A_C}^2=(\ell-2)\frac{R}{T}+\frac{2(\ell-1)}R.
\end{equation}
On $\Lambda=C\cap\Sph^{2\ell-1}$, put $t=T|_\Lambda$. Then
\begin{gather}\label{bh:link-table}
 \norm{A_\Lambda}^2=\frac{\ell-2}{t}+2(\ell-1),\qquad
 \norm{\nabla_\Lambda t}^2=4t(1-4t),\\
 \Delta_\Lambda t=2(\ell-1)-8\ell t.\nonumber
\end{gather}
\end{lemma}
\begin{proof}
Put $W=\norm{\nabla F}^2$ and $\nu=\nabla F/\sqrt W$.
The Euclidean Bochner identity gives
$\norm{\nabla^2F}^2=\Delta W/2-\Gamma(F,\Delta F)$.
On $F=0$, \eqref{bh:reduced-table} yields
\begin{align*}
 W&=4bRT,\\
 \norm{\nabla^2F}^2&=4(a^2+b\ell)R^2+8b(\ell+4)T,\\
 \norm{\nabla\sqrt W}^2&=4[(a^2+b)R^2+5bT],\\
 \nabla^2F(\nu,\nu)&=\Delta F=2aR.
\end{align*}
Now use
\[
 \norm{A_C}^2=
 \frac{\norm{\nabla^2F}^2-2\norm{\nabla\sqrt W}^2
                     +(\nabla^2F(\nu,\nu))^2}{W}
\]
to obtain \eqref{bh:AC}. The normal Hessian of any polynomial $V$ can also be computed
from the closure table, using
\[
 \Hess V(\nu,\nu)=\frac{
 \Gamma(F,\Gamma(F,V))-\tfrac12\Gamma(V,W)}{W}.
\]
On $F=0$, this gives
\[
 \Gamma(T,F)=0,\qquad
 \Hess T(\nu,\nu)=2R-\frac{4T}{R},\qquad
 \Delta_C T=2(\ell-1)R+\frac{4T}{R}.
\]
Subtracting the radial contribution of the degree-four function $T$
on the $(2\ell-1)$-dimensional cone gives the last identity in
\eqref{bh:link-table}. Subtracting the squared radial derivative from
$\Gamma(T,T)=4RT$ gives the gradient identity.
\end{proof}

For $\ell=8$, the positive function $\phi=t^{-1}$ satisfies
\begin{equation}\label{bh:eigenfunction}
 (\Delta_\Lambda+\norm{A_\Lambda}^2)\phi=46\phi.
\end{equation}
To apply this identity, we cut off near the singular set of the link.
Its two singular strata have dimension seven and
transverse dimension seven in $\Lambda$. The local quadratic-cone
description after \eqref{bh:transverse} implies
$t\asymp\rho^2$ and transverse volume density comparable to
$\rho^6\,d\rho$, where $\rho$ is distance to the singular stratum.
Thus $\phi\asymp\rho^{-2}$, its gradient is $O(\rho^{-3})$, and
$\phi$, its gradient and the potential term have the required
finite integrals.

Choose cutoffs $\eta_\varepsilon$ vanishing for $\rho<\varepsilon$
and equal to one for $\rho>2\varepsilon$, with gradient bounded by
$O(\varepsilon^{-1})$. Then
\[
 \int_\Lambda\phi^2\norm{\nabla\eta_\varepsilon}^2
 =O(\varepsilon^{-4}\varepsilon^{-2}\varepsilon^7)
 =O(\varepsilon)\longrightarrow0.
\]
Multiplying \eqref{bh:eigenfunction} by
$\eta_\varepsilon^2\phi$ and integrating by parts on the regular
part proves that the angular Rayleigh quotients of
$\eta_\varepsilon\phi$ tend to $-46$. The cone has dimension fifteen,
so \eqref{eq:exact-Hardy-gap} gives compactly supported regular
variations whose scale-invariant quotients tend to
\begin{equation}\label{bh:negative-margin}
 \frac{(15-2)^2}{4}-46=-\frac{15}{4}<0.
\end{equation}
This proves instability. Together with \eqref{bh:Lawson-margin} and
the evenness of $\ell$, it completes Theorem~\ref{bh:main}.

\subsubsection{The Hopf pairing and its balanced section}
\label{sec:Hopf-pairing}

The Hopf pairing realizes weights $(1,8,8)$ and has two linear
singular strata. Its balanced section also gives an intrinsic
comparison inside a Simons cone.

Let \(h:\R^{16}\to\R^9\) be the quadratic extension of the standard
octonionic Hopf map \(\Sph^{15}\to\Sph^8\), normalized by
\begin{equation}\label{eq:Hopf-normalization}
 |h(x)|^2=|x|^4,\qquad
 Dh(x)Dh(x)^{\mathsf T}=4|x|^2\Id_9,\qquad \Delta h=0.
\end{equation}
For polarization of the spherical norm identity and the Hopf construction
from norm-preserving bilinear multiplication, see
\cite[Chapter~15]{Shapiro2000}.
The coordinates are \(h_i(x)=\langle A_ix,x\rangle\) for a nine-matrix Clifford
system on \(\R^{16}\); see \cite[Theorem~A and Section~1.2]{Radeschi14}.
On \(\R^{16}\oplus\R^{16}\), set
\begin{equation}\label{eq:Hopf-pairing-definitions}
 u=|x|^2,\quad v=|y|^2,\quad R=u+v,\quad V=u-v,\quad T=uv,
 \qquad F(x,y)=\langle h(x),h(y)\rangle.
\end{equation}
In particular, \(|F|\leq T\), and \(F\) has bidegree \((2,2)\).

\begin{proposition}[The pairing closure]\label{prop:Hopf-closure}
The algebra \(\cA=\R[u,v,F]=\R[R,V,F]\) is a graded Laplacian
algebra of spherical quotient dimension two. Its closure table, in
the ordered generators \((u,v,F)\), is
\begin{equation}\label{eq:Hopf-closure}
 (\Delta u,\Delta v,\Delta F)=(32,32,0),\qquad
 \bigl(\Gamma(I_i,I_j)\bigr)=
 \begin{pmatrix}
  4u&0&4F\\
  0&4v&4F\\
  4F&4F&4uv(u+v)
 \end{pmatrix}.
\end{equation}
The quartic zero set \(C=\{F=0\}\) is minimal on its regular part,
and its geometric singular set is
\begin{equation}\label{eq:Hopf-singular-set}
 Z=(\R^{16}\times\{0\})\cup(\{0\}\times\R^{16}).
\end{equation}
\end{proposition}

\begin{proof}
Equation \eqref{eq:Hopf-normalization} gives
\[
 |\nabla_xF|^2=4uv^2,\qquad
 |\nabla_yF|^2=4u^2v,\qquad \Delta F=0.
\]
Euler's identity in each block gives the remaining entries of
\eqref{eq:Hopf-closure}, and \cref{prop:MR-generator-test}
proves Laplacian closure. Surjectivity of the Hopf map gives
the full image
\[
 \{(u,v,f):u,v\geq0,\ |f|\leq uv\}.
\]
Its interior proves algebraic independence and spherical quotient
dimension two. Since $uv$ varies at fixed $R$ and $F=0$,
$\Gamma(F,F)=4Ruv\notin\R[R,F]$.

Since \(\Gamma(F,R)=8F\) and \(\Gamma(F,T)=4RF\), we obtain
\[
 \cL(F)=-\tfrac12\Gamma(F,4RT)=-8(R^2+2T)F.
\]
This proves regular minimality and identifies the critical set as
\(Z\). At \((0,y_0)\), \(y_0\ne0\), the transverse quadratic term
is \(\langle A(h(y_0))x,x\rangle\), where
\(A(a)=\sum_i a_iA_i\). Its eigenvalues are
\(\pm|y_0|^2\), each with multiplicity eight. Hence the tangent
cone is the product of a quadratic cone of signature \((8,8)\)
with \(\R^{16}\), and the point is geometrically singular.
The other axis is identical; the vertex is singular as well.
\end{proof}

\begin{remark}[A second realization of the pairing table]
The Clifford construction in \cref{sec:Clifford-composition-background}
also realizes this table. Put
\(Q=u^2+v^2-2F\). For a nine-matrix Clifford system
\(P_0,\ldots,P_8\) on \(\R^{32}\), the generators
\[
 \widehat R=|z|^2,\qquad
 \widehat V=\langle P_0z,z\rangle,\qquad
 \widehat Q=\sum_{i=0}^8\langle P_iz,z\rangle^2
\]
have the same gradient and Laplacian identities as \((R,V,Q)\).
Both images are
\(\{(R,V,Q):R\geq0,\ V^2\leq Q\leq R^2\}\).
The inverse relations
\(u=(R+V)/2\), \(v=(R-V)/2\), and \(F=(R^2+V^2-2Q)/4\)
therefore identify their graded differential algebras.
\end{remark}

\begin{theorem}[A polynomial subcalibration of the pairing cone]
\label{thm:Hopf-pairing}
The potential \(G=TF=|x|^2|y|^2F\) has degree eight and
subcalibrates \(C=\{F=0\}\). Its oriented boundary current
\(\partial\llbracket\{F<0\}\rrbracket\) is strictly area minimizing
in \(\R^{32}\). Moreover,
\begin{equation}\label{eq:Hopf-Hardy-bound}
 Q_C(\varphi)\geq\frac{449}{4}\int_C r^{-2}\varphi^2,
 \qquad \varphi\in C_c^\infty(C_{\reg}),\quad r=\sqrt R.
\end{equation}
\end{theorem}

\begin{proof}
The three blocks
\[
 E_0=(u-v)^2,\qquad E_1=2(T+F),\qquad E_2=2(T-F)
\]
have weights $(1,8,8)$. The centered generator $F/2$ gives
$h=16$, $r_0=1$, and $c=1/2$ in
\cref{thm:quartic-table-existence}. Its balanced specialization
\cref{prop:balanced-quartic-multiplier} gives $G=TF$ and
\begin{align}
 \Delta G&=40RF,\qquad
 |\nabla G|^2=4RT(T^2+3F^2),\label{eq:Hopf-gradient}\\
 \cL(G)&=32TF\bigl[T^2(2R^2-T)
                    +3F^2(4R^2-T)\bigr].\label{eq:Hopf-defect}
\end{align}
The Hardy constant in \eqref{eq:bihomogeneous-Hardy} is $449/4$
at $(N,h)=(32,16)$.
\end{proof}

\begin{remark}[Relation with the OT--FKM and Riedler constructions]
\label{rem:Hopf-FKM-Riedler}
The matrices \(B_i=\operatorname{diag}(A_i,-A_i)\) form a Clifford
system on \(\R^{32}\). Its OT--FKM polynomial
\cite[Section~1.3]{Radeschi14} is
\[
 \Psi=R^2-2\sum_{i=1}^9\langle B_i(x,y),(x,y)\rangle^2
     =4F-V^2.
\]
It has multiplicities \((8,7)\). Applying the centering formula
\eqref{eq:iso-F-z} to \(-\Psi\), with ordered multiplicities
\((7,8)\), shows that its minimal regular hypercone is
\(\{\Psi+R^2/15=0\}\). Nonvertex singularities distinguish $C$ from every cone over a
regular isoparametric hypersurface. They also distinguish it from
Riedler's family: its singular dimension is $N-2k$, so congruence
would require $N=32$, $2k=16$, but sixteen Clifford matrices
require ambient dimension at least $256$
\cite[Remark~5.2.6(1) and Table~4.1]{RiedlerThesis25}. Although
\(F=\langle\cdot,\cdot\rangle\circ(h,h)\), the product map
\(\Phi(x,y)=(h(x),h(y))\) satisfies
\[
 D\Phi D\Phi^{\mathsf T}=4\operatorname{diag}(u\Id_9,v\Id_9).
\]
It is not horizontally conformal on \(u\ne v\), so the pullback
criterion in \cref{prop:pullback-subcalibration} does not apply
to this map on \(\R^{32}\).
\end{remark}

\paragraph{The balanced section inside a Simons cone.}
\label{sec:Hopf-intrinsic-extension}

For the pairing polynomial above, the ambient space for comparison
is now itself a minimal cone. Let
\[
 \mathcal B=\{V=0\}=\{|x|=|y|\},\qquad
 D=\{V=F=0\},\qquad E_{\mathcal B}=\mathcal B\cap\{F<0\}.
\]
The Simons cone \(\mathcal B\) is a singular space of dimension
\(31\), smooth away from the origin. We use its induced metric
and the orientation determined by \(\nabla V\).
Intrinsic perimeter is variation in the smooth part, extended across
the vertex by the cutoff below. Write it as \(\Per_{\mathcal B}\), and set
\(D_a=\partial\llbracket E_{\mathcal B}\rrbracket
\mathbin{\llcorner}B_a(0)\).

\begin{theorem}[The pairing section of a Simons cone]
\label{thm:Hopf-Simons-section}
The cone \(D\) has dimension \(30\), is smooth away from the
origin, and is minimal in \(\R^{32}\). The intrinsic potential
\(G_{\mathcal B}=RF\) subcalibrates it in \(\mathcal B\).
For every \(a>0\), the integral current \(D_a\) minimizes mass
among integral \(30\)-currents supported in
\(\mathcal B\cap\overline B_a(0)\) with boundary \(\partial D_a\).
Thus \(D_a\) solves the Plateau problem in the singular ambient
space \(\mathcal B\), with its Euclidean induced area.
\end{theorem}

\begin{proof}
On \(\mathcal B\setminus\{0\}\), the gradients of \(V\) and
\(F\) are nonzero and orthogonal by \eqref{eq:Hopf-closure}.
Hence \(D\setminus\{0\}\) is smooth of codimension two.
The unit normal to \(\mathcal B\) is
\(n_{\mathcal B}=(x,-y)/\sqrt R\). The bidegree of \(F\) gives
\(\Hess F((x,-y),(x,-y))=-4F\). Since \(\mathcal B\) is minimal,
the induced operators satisfy
\begin{equation}\label{eq:Hopf-intrinsic-closure}
 \begin{gathered}
 \Delta_{\mathcal B}R=62,\qquad
 \Delta_{\mathcal B}F=4F/R,\\
 \Gamma_{\mathcal B}(R,R)=4R,\qquad
 \Gamma_{\mathcal B}(R,F)=8F,\qquad
 \Gamma_{\mathcal B}(F,F)=R^3.
 \end{gathered}
\end{equation}
Here \(\nabla F\) is tangent to \(\mathcal B\), and \(T=R^2/4\).
The same chain rule as in \cref{thm:finite-reduction} now yields
\begin{equation}\label{eq:Hopf-intrinsic-defect}
 \begin{gathered}
 \Delta_{\mathcal B}(RF)=82F,\qquad
 |\nabla_{\mathcal B}(RF)|^2=R^5+20RF^2,\\
 \cL_{\mathcal B}(RF)=16RF(2R^4+85F^2),
 \end{gathered}
\end{equation}
where \(\cL_{\mathcal B}\) is \eqref{eq:L-definition} for the
induced metric. The normalized gradient
\[
 X=\frac{\nabla_{\mathcal B}(RF)}{|\nabla_{\mathcal B}(RF)|}
\]
is smooth on \(\mathcal B\setminus\{0\}\), equals the exterior
normal along \(D\setminus\{0\}\), and its divergence has the
strict sign of \(F\) off \(D\).

We apply the proof of \cref{thm:removable-subcalibration} in the
smooth part of \(\mathcal B\). Only the vertex requires a cutoff.
For a radial cutoff \(\eta_\eps\), zero on \(B_\eps(0)\) and
one outside \(B_{2\eps}(0)\),
\[
 \int_{\mathcal B\cap B_{2\eps}(0)}
       |\nabla_{\mathcal B}\eta_\eps|\,d\cH^{31}
 =O(\eps^{30}).
\]
Hence extending a bounded characteristic function across the vertex
adds no perimeter mass. The estimate $|\diver_{\mathcal B}X|\le c/r$
from homogeneity makes the divergence error vanish as well. For every relatively compact open set \(U\subset\mathcal B\)
and every locally finite-perimeter set \(H\subset\mathcal B\) with
\(H\triangle E_{\mathcal B}\Subset U\),
\begin{equation}\label{eq:Hopf-intrinsic-comparison}
 \Per_{\mathcal B}(H;U)-\Per_{\mathcal B}(E_{\mathcal B};U)
 \geq\int_{H\triangle E_{\mathcal B}}
       |\diver_{\mathcal B}X|\,d\cH^{31}\geq0.
\end{equation}
The set \(E_{\mathcal B}\) has locally finite perimeter because
its boundary is a cone with a smooth compact link.

For the asserted integral-current comparison, cone a competitor
minus \(D_a\) to the origin inside \(\mathcal B\cap\overline B_a(0)\).
This gives an integral \(31\)-dimensional filling. On the oriented smooth part of \(\mathcal B\), write this filling as
\(\llbracket k\rrbracket_{\mathcal B}\), with integer-valued
locally BV multiplicity, and set
\(v=\one_{E_{\mathcal B}}+k\) and \(H=\{v\geq1\}\).
The integer-level decomposition used in
\cref{thm:strict-minimality-flux} gives
\(\lvert D\one_H\rvert\leq|Dv|\) on the smooth part.
Writing \(T\) for the competitor, for \(b>a\) we obtain
\[
 0\leq \Per_{\mathcal B}(H;B_b)-
          \Per_{\mathcal B}(E_{\mathcal B};B_b)
 \leq \mathbf M(T)-\mathbf M(D_a).
\]
The inequalities use \eqref{eq:Hopf-intrinsic-comparison} and
cancellation of the common exterior mass; $D$ has no mass on a sphere.
The representation and integer-level decomposition hold in smooth
charts. The vertex cutoff extends the characteristic function,
and no positive-dimensional integral current is supported at the
vertex. Thus the comparison includes competitors meeting it.

Finally, the divergence in \eqref{eq:Hopf-intrinsic-defect} vanishes
on \(D\), so \(D\) is minimal in \(\mathcal B\).
For its unit normal \(\nu\) within \(\mathcal B\), the two block
lengths satisfy \(|\nu_x|=|\nu_y|\), by the partial gradient
identities above and \(u=v\). Hence
\(\mathrm{II}_{\mathcal B}(\nu,\nu)=0\). Since
\(\operatorname{tr}_{T\mathcal B}\mathrm{II}_{\mathcal B}=0\),
its trace on \(TD\) is also zero. The second fundamental form
composition formula proves Euclidean minimality of \(D\setminus\{0\}\).
The smooth compact link gives
\(\int_{D\cap B_{2\eps}}|\nabla_D\eta_\eps|\,d\cH^{30}
=O(\eps^{29})\to0\), so the Euclidean first variation also vanishes
across the vertex.
\end{proof}

The Euclidean area-minimizing cones of codimension two in
\cite[Theorem~B]{CuiFKM25} are also area-minimizing hypersurfaces within
the Simons cones containing them. Indeed, let \(K\subset\mathcal S\)
be such a cone and put \(K_a=K\mathbin{\llcorner}B_a(0)\).
Every integral current \(T\) supported in
\(\mathcal S\cap\overline B_a(0)\) with
\(\partial T=\partial K_a\) is also an admissible Euclidean
competitor, so \(\mathbf M(K_a)\leq\mathbf M(T)\).
The induced area on \(\mathcal S\) is the Euclidean area.




\subsubsection{Riedler's Clifford--Simons quartics}

The construction and area-minimizing range
\(N\geq32\), \(k\geq4\) are due to Riedler
\cite[Section~5.2 and Theorem~5.2.5]{RiedlerThesis25}.
The common quartic theorem gives polynomial potentials; the
degree-nine potential below improves their Hardy bound. The
sixteen-dimensional endpoint is a splitting cone.

Let \(A_1,\ldots,A_{2k}\) be a symmetric Clifford system on
\(\R^N\), and write
\(P_i(x)=\langle A_i x,x\rangle\) and
\(Q=\sum_{i=1}^{2k}P_i^2\). Define
\[
 F=\sum_{i=1}^{k}P_i^2-\sum_{i=k+1}^{2k}P_i^2,
 \qquad C^4_{N,k}=\{x\in\R^N:F(x)=0\}.
\]
This particular Clifford realization requires
$\delta(2k-1)\mid N/2$ \cite[Section~3.5]{FKM81}.
Here $F=f\circ P$, where
$f(y)=\sum_{i\leq k}y_i^2-\sum_{i>k}y_i^2$, so the composition
is the pullback of a Simons cone. The unequal-target example
\cref{ex:lawson-pullback} extends this connection with the transfer
theorem of \cref{sec:quartic-application}.
At the comparison endpoint \(k=1\),
\(F=(P_1-P_2)(P_1+P_2)\) is reducible and harmonic;
\cref{prop:minimizing-irreducibility} excludes area minimization.
For the minimizing parameters below, the same proposition gives
irreducibility.
The chain rule gives the identities recorded in
\cite[Lemma~5.2.11]{RiedlerThesis25}:
\begin{equation}\label{eq:quartic-ambient-closure}
 \bigl(\Gamma(I_i,I_j)\bigr)_{I=(R,Q,F)}=
 \begin{pmatrix}
 4R&8Q&8F\\
 8Q&16RQ&16RF\\
 8F&16RF&16RQ
 \end{pmatrix},\qquad
 \Delta I=(2N,16kR,0).
\end{equation}
The determinant is \(1024R(R^2-Q)(Q^2-F^2)\), so the spherical
quotient has dimension two wherever \(0<Q<R^2\) and \(|F|<Q\).
For a concrete example, take the quadratic octonionic Hopf map
\(\widetilde P=(P_1,\ldots,P_9):\R^{16}\to\R^9\), normalized by
\(|\widetilde P|^2=R^2\), and retain its first eight components.
Then \(k=4\) and \(Q=R^2-P_9^2\). Choose an image with
\(P_9\neq0\) and nonzero components in both four-dimensional blocks.
This gives \(0<Q<R^2\) and \(|F|<Q\). Taking two copies of this
eight-matrix Clifford system gives \(N=32\), \(k=4\), in the
minimizing range below; the same point, placed in one copy, retains
the strict inequalities.
When \(Q\not\equiv R^2\), these generators correspond to
\eqref{eq:orthogonal-quartic-data} via
\[
 (E_0,E_1,E_2)=\left(R^2-Q,\frac{Q+F}{2},\frac{Q-F}{2}\right),
 \qquad (r_0,r_1,r_2)=\left(\frac N2-2k+1,k,k\right).
\]
For $N\ge32$, $k\ge4$, normalize $T=Q/4$, $F_{\rm gen}=F/8$.
Lemma~\ref{lem:quartic-table-geometry} gives regular minimality,
the singular set $\{Q=0\}$, and full support. If $k=4$, then
$r_0\ge9$ and \cref{thm:quartic-table-existence} gives $G=RQF$;
if $k\ge5$, \cref{prop:balanced-quartic-multiplier} gives $G=QF$.
Thus the boundary current is strictly area minimizing in this range,
with the Hardy bound \eqref{eq:bihomogeneous-Hardy}. Positive
constant rescalings do not affect normalized gradients.

The defining polynomial itself satisfies
\[
 |\nabla F|^2=16RQ,\qquad \Delta F=0,\qquad
 \cL(F)=-64F(Q+2R^2).
\]
The last identity follows from
$\cL(F)=-\Gamma(F,\Gamma(F,F))/2$ and
\eqref{eq:quartic-ambient-closure}.

\begin{example}[Riedler's subcalibration potential]
\label{ex:Riedler-correct-power}
For \(N\geq32\) and \(k\geq4\), Riedler's thesis uses
\cite[Section~5.2.2, Lemmas~5.2.12--5.2.13]{RiedlerThesis25}
\begin{equation}\label{eq:Riedler-correct-potential}
 G=RQ^{3/4}F=16^{-3/4}R^{1/4}S^{3/4}F,
 \qquad S=|\nabla F|^2=16RQ.
\end{equation}
This potential has degree nine. The parameters in
\cref{prop:harmonic-power-Jacobi} are
\((e,\sigma,\gamma)=(4,1/2,3/4)\).
Using \eqref{eq:quartic-ambient-closure}, we obtain the boundary Jacobi coefficient
\begin{equation}\label{eq:Riedler-positive-boundary}
 \beta_G=(2N-41)+(12k-47)\frac{R^2}{Q}>0.
\end{equation}
Since \(Q\leq R^2\), \cref{thm:ground-state-gap} gives
\[
 Q_C(u)\geq
 \left[\left(9-\frac{N-1}{2}\right)^2+2N+12k-88\right]
 \int_Cr^{-2}u^2.
\]
This lower bound exceeds the general quartic bound
\eqref{eq:bihomogeneous-Hardy}, with $h=2k$, by
\[
 \left[\frac{(N-19)^2}{4}+2N+12k-88\right]
 -\left[\frac{(N-17)^2}{4}+16k-72\right]
 =N+2-4k=2r_0>0.
\]
Thus Riedler's degree-nine potential gives the stronger Hardy bound.
\end{example}

\begin{corollary}[Instability of Riedler's sixteen-dimensional quartic]
\label{thm:Riedler-unstable}
The cone $C^4_{16,4}$ is unstable on its regular part and is not area
minimizing. More precisely, there are compactly supported variations on the regular
part with scale-invariant quotients of the second variation tending to
$-15/4$.
\end{corollary}

\begin{proof}
Write $F_R$ for the defining quartic of $C^4_{16,4}$ and let
$A_1,\ldots,A_8$ be its symmetric Clifford matrices. The volume
matrix $A_9=A_1\cdots A_8$ satisfies
\[
 A_9^T=A_9,\qquad A_9^2=I,\qquad A_9A_i=-A_iA_9.
\]
Its two eigenspaces have dimension eight, since $A_1$ exchanges
them isometrically. Choosing compatible orthonormal bases gives
coordinates $(x,y)\in\R^8\oplus\R^8$ in which
\[
 A_{i+1}=\begin{pmatrix}0&B_i\\B_i^T&0\end{pmatrix}
 \quad(0\leq i\leq7),\qquad B_0=I.
\]
The matrices $B_1,\ldots,B_7$ are skew-symmetric and satisfy
$B_iB_j+B_jB_i=-2\delta_{ij}I$. Thus the eight vectors $B_i y$
$(0\leq i\leq7)$ are pairwise orthogonal, all of length $|y|$, and
\begin{gather*}
 \sum_{i=0}^7\langle x,B_i y\rangle^2=|x|^2|y|^2,
 \qquad Q=4|x|^2|y|^2,\\
 \frac{F_R}{8}
 =\sum_{i=0}^3\langle x,B_i y\rangle^2
       -\frac12|x|^2|y|^2=F_{\rm split}.
\end{gather*}
Consequently $C^4_{16,4}$ is orthogonally congruent to a splitting
cone \eqref{bh:centered} with $(\ell,m)=(8,4)$.
Its Clifford representation is the balanced one: the symmetric
involution $B_1B_2B_3$ anticommutes with $B_4$, so its two
eigenspaces have dimension four.

In the original variables, $\Lambda=C^4_{16,4}\cap\Sph^{15}$ and
$q=Q|_\Lambda=4t$, so \eqref{bh:link-table} reads
\begin{equation}\label{eq:Riedler-16-link}
 |A_\Lambda|^2=\frac{24}{q}+14,\qquad
 |\nabla_\Lambda q|^2=16q(1-q),\qquad
 \Delta_\Lambda q=56-64q.
\end{equation}
The $\ell=8$ argument above transfers under this orthogonal
identification, and \eqref{bh:negative-margin} gives the asserted
quotients and instability.
\end{proof}

This is the case left undecided in
\cite[Remark~5.2.6(3)]{RiedlerThesis25}. Together with
Theorem~5.2.5 and the exclusions in Remark~5.2.6(1)--(2)
of that thesis, it gives
\[
 C^4_{N,k}\text{ is area minimizing}
 \quad\Longleftrightarrow\quad N\geq32\text{ and }k\geq4
\]
among the admissible Clifford--Simons quartics.

